\chapter{TEMP Funkcje wielu zmiennych}

\section{Ekstrema funkcji wielu zmiennych.}

W analizie jednej zmiennej, aby znaleźć ekstrema, przyrównujemy pierwszą pochodną do zera i sprawdzamy znak drugiej pochodnej. W przypadku funkcji $f: \mathbb{R}^{n}\to \mathbb{R}$ sytuacja jest analogiczna, choć narzędzia są trochę inne -- zamiast pochodnej jest \textbf{gradient}, a zamiast drugiej pochodnej -- macierz Hassego. 

Zacznijmy od ścisłej definicji minimum lokalnego. Zwróćmy uwagę na warunek: wartość funkcji w punkcie $x_0$ musi być najmniejsza w pewnym otoczeniu tego punktu.

\begin{definition}[Lokalne minimum (maksimum)]
	Załóżmy, że funkcja ciągła $f$ o wartościach rzeczywistych jest określona w pewnym podzbiorze $\mathbb{R}^{n}$. Mówimy, że punkt $x_{0}$ jest \textbf{lokalnym minimum} funkcji $f$ jeśli dla pewnej liczby $\delta > 0$ mamy $f\left( x \right) \ge f\left( x_0 \right) $ dla $\|x-x_0\| < \delta$.

Analogicznie definiujemy \textbf{lokalne maksimum} 
\end{definition}  
